\section{Technical supplement: wall proofs, regressions, and finite evidence}

\subsection{An exact bounded DNA chamber atlas}
\label{sec:finite-atlas}

We retain the original thirty-row regression result alongside the complete
order-eight classification in \Cref{thm:complete-atlas}. It uses the DNA family
\eqref{eq:dna-weights}, the complete catalog of \Cref{def:decision-catalog},
exact algebraic sign evaluation, and the deterministic selector of
\Cref{thm:imported-selector}.  The range is exactly thirty rows: every
interval and critical point for $k=4,6$, and four named weights for $k=8$.

The landmarks are
\[
  \begin{array}{c|c}
  \role{rat2} & 2\\
  \role{gold5} & (7+3\sqrt5)/2\\
  \role{quad13} & \sqrt{13}\\
  \role{cubic2} & \sqrt[3]{2}.
  \end{array}
\]

\begin{proposition}[Exact small DNA chamber atlas]
\label{prop:finite-atlas}
Under the conventions above:
\begin{enumerate}
  \item The complete $k=4$ diagram has three cells and two terminal/residual
        classes.  Its sole critical ratio is $\rho=3$; the class for
        $1<\rho\le3$ differs in sixteen selected words from the class for
        $\rho>3$.  Both have longest residual path $17$.
  \item The complete $k=6$ diagram has the eleven ordered critical ratios
  \[
    \frac75,\frac32,\frac53,2,\frac73,3,4,5,7,9,11,
  \]
  hence twenty-three point/interval cells and seventeen terminal/residual
  classes.  Its longest-path phase is
  \[
    \begin{array}{c|c}
    1<\rho\le7/3 & 63\\
    7/3<\rho\le3 & 65\\
    3<\rho<4 & 59\\
    \rho\ge4 & 57.
    \end{array}
  \]
  Except at $\rho=3$, every cell has $144$ lower zeros and $220$ physical
  firings; at $\rho=3$ the counts are $176$ and $212$.
  \item At $k=8$, the four declared landmarks give four distinct terminal
        sets, midpoint ideals, and residual digraphs, although all four have
        $256$ zero-embedding words, $70$ zero PCRs, $576$ lower-zero states,
        $3952$ firings, and minimum cardinality $8230$.  Their longest paths
        are respectively $147,143,141,161$ in the landmark order above.
  \item Every unequal pair in the frozen landmark comparisons differs outside
        zero-PCR tie representatives.  No equal complete covector produced
        two different terminal payloads.
\end{enumerate}
\end{proposition}

\begin{proof}[Computer-assisted proof]
All thirty cells were generated with exact arithmetic over rational and
represented real algebraic numbers.  For each row, the producer enumerated all labelled words,
constructed the seed and midpoint ideal, formed the direct terminal, and
checked literal reverse-complement invariance, one representative per PCR,
and acyclicity.  It then computed the labelled residual digraph and longest
path.

An independent checker reconstructed the complete decision covector,
terminal, legal firing replay, residual arcs, topological order, path
histogram, and all pairwise comparisons.  Two sequential producer runs and
normal/optimized checker runs had identical canonical scientific payloads.
The assurance layer rejected twenty-nine typed semantic mutations, including
reversed gradients, omitted decision forms, floating zero tests, changed tie
rules, incomplete cell sets, and altered path statistics.  The canonical
producer and validator payload hashes are recorded in
\Cref{app:finite-evidence}.
\end{proof}

This regression proposition alone does not prove the complete $k=8$
diagram; \Cref{thm:complete-atlas} supplies that extension.
The incompatibility of optimal weights is resolved in
\Cref{prop:family-optima}.  In fact the $k=6$ landmark ranking favors \role{gold5},
whereas among the four $k=8$ controls \role{quad13} is shorter.  Equal zero
counts can coexist with thousands of selected-word changes, showing why
zero-sector extremality and residual geometry are distinct objectives.

\subsection{Wall geometry and the complete DNA atlas}
\label{sec:walls-atlas}
Write $\theta=\RC$ on lower vertices, $J=\RC$ on arcs, and
$\alpha=2\pi/k$. We use the gradient convention \eqref{eq:gradient}.
The basic identity, also used in \Cref{ext:thm:paths}, is
\begin{equation}
 r_s(v)-r_s(u)=\tan(\pi/k)\bigl(\lambda_s(u)+\lambda_s(v)\bigr)
 \quad (u\to v).
 \label{eq:wall-increment}
\end{equation}

\subsubsection{Regular midpoint walls}
Consider an admissible affine weight path $s_t=s_0+t\dot s$. Call its
wall at zero \emph{regular} if every lambda form either vanishes
identically on this path or is nonzero at zero. Restrict to a sufficiently
small neighborhood with no other roots. Let
\[
 S=\{u:\lambda_{s_0}(u)>0,\ r_{s_0}(u)=0,\ \dot r(u)\ne0\}.
\]
Subscripts $-,0,+$ denote the two exact side germs and the point.

\begin{theorem}[Regular wall locality]\label{thm:regular-wall}
The set $S$ is RC-stable and no two of its vertices are adjacent.
The seed is constant and
\[
 m_+-m_-=\nabla f,\qquad
 f(u)=-\sgn\dot r(u)\ \ (u\in S),\quad f=0\ \text{otherwise}.
\]
Thus $|M_+\triangle M_-|=2q|S|$. Between any two of the three
germs the same formula holds using only the vertices whose midpoint
indicator changes. Their three lower residual lengths lie in an
interval of width at most two.
\end{theorem}
\begin{proof}
The two-family formula makes the seed constant. Only vertices of $S$
can change $a_s$, with the stated signs. Identity
\eqref{eq:wall-increment} rules out adjacent vertices of $S$; it also
rules out loops there. The incident stars are disjoint. On each star
both endpoint indicators are Boolean, so the change exchanges all
$q$ incoming and $q$ outgoing arcs. In each residual a changing
vertex is a source or a sink. A path meets these vertices only at
its endpoints. Deleting them leaves a path in the same induced
residual outside $S$; hence every length lies between the longest
path there and that length plus two. Finally
$\lambda(\theta u)=\lambda(u)$, $r(\theta u)=-r(u)$ and
$a(u)+a(\theta u)=1$ on $S$, including the frozen point tie.
This proves RC stability and all point comparisons.
\end{proof}
The bound two is attained at $k=8$, $\rho=4\sqrt2-1$, where the
lengths are $(145,143,143)$. It does not apply to a mixed wall.

\subsubsection{Singular vertices and mixed walls}
Put $Z=\{z:\lambda_{s_0}(z)=0\}$, including identically zero forms.
For a letter $c$ define
\[
 x_c(z)=\Re P_{s_0}(zc).
\]
If $y_c$ is the tail of $cz$ and $v_c$ the head of $zc$, rotation gives
\[
 \lambda(y_c)=x_c(z)\sin\alpha,\qquad
 \lambda(v_c)=-x_c(z)\sin\alpha.                     \tag{W1}
\]

\begin{lemma}[Singular subgraph]\label{lem:singular-subgraph}
The original graph induced on $Z$ is a disjoint union of isolated
vertices and zero-embedding PCR cycles. Each cycle has length dividing
$k$, and its midpoint potential is constant, equal to $-W_s(w)$.
A zero PCR with nonzero embedding derivative along the weight path
has period exactly $k$.
\end{lemma}
\begin{proof}
Since $x_c(z)$ is a fixed real number plus $s_0(c)$, injectivity
allows at most one letter with $x_c(z)=0$. Equation (W1) shows that
this letter gives both the unique incoming and the unique outgoing
internal arc. Thus internal indegree and outdegree are either both
zero or both one. Rotation of a zero embedding supplies the
successor, so the cycles are precisely the zero PCRs, without tails.
Their total letter weight is rotation invariant, proving the
potential formula. A word of proper period $d\mid k$ has embedding
zero for every weight, by the geometric sum across repeats; its
derivative cannot be nonzero.
\end{proof}

For any two germs $\sigma,\tau\in\{-,0,+\}$, put
$\Delta a=a_\tau-a_\sigma$ and
$\Delta h=\ind_{\lambda_\tau>0}-\ind_{\lambda_\sigma>0}$. For
$x_c(z)\ne0$ define
\[
 \gamma_c(z)=\Delta a(z)-\Delta h(z)\ind_{x_c(z)<0}.
\]
Let $e_\tau(C)$ be the unique arc of a singular PCR $C$ selected
by terminal $\tau$. Keep the definition of $S$ from the regular case.

\begin{theorem}[Disjoint mixed-wall formula]\label{thm:mixed-wall}
One has the exact cochain identity
\begin{align}
 m_\tau-m_\sigma={}&\nabla(\ind_S\Delta a)
 +\sum_{\substack{z\in Z,\ c\in\A\\x_c(z)\ne0}}
       \gamma_c(z)(\ind_{zc}-\ind_{cz}) \notag\\
 &+\sum_{C\subset G[Z]}(\ind_{e_\tau(C)}-\ind_{e_\sigma(C)}).
 \label{eq:mixed-wall}
\end{align}
All nonzero summands have disjoint supports, except for the
incoming/outgoing arcs already grouped into one star. In particular,
\[
 |M_\tau\triangle M_\sigma|=2q\,s_{\sigma\tau}
                       +2e_{\sigma\tau}+2c_{\sigma\tau},
\]
where $s$ counts changed midpoint vertices in $S$,
$e=\sum_{z,c}|\gamma_c(z)|$, and $c$ counts singular PCRs whose
selected representative changes.
\end{theorem}
\begin{proof}
Outside $Z$, only midpoint indicators at $S$ change. On $zc$ with
$x_c(z)\ne0$, the seed is $(1-h_\tau(z))h(v_c)$. After separating
a possible midpoint change at $v_c$, its difference is
$\Delta a(z)-\Delta h(z)h(v_c)=\gamma_c(z)$ by (W1).
The companion $cz$ has the negative difference. The inequalities
$0\le a_\tau\le h_\tau\le1$ give $\gamma_c(z)\in\{-1,0,1\}$.

For the boundary check, if $u\in S$ and $u\to z$ with $z\in Z$,
\eqref{eq:wall-increment} gives $r(z)>0$, so $a_\tau(z)=0$
in all germs and the singular contribution on that arc vanishes.
If $z\to u$, then $r(z)<0$, so $a_\tau(z)=h_\tau(z)$ and the
same conclusion holds. Thus nonzero singular contributions do not
overlap midpoint stars. Each external arc has only one endpoint
in $Z$; a companion coincidence is a homopolymer loop and has
$x_c(z)=0$, hence is excluded. Internal arcs belong to the disjoint
PCRs of \Cref{lem:singular-subgraph}, each with one selected arc.
All other arcs are unchanged. Summing absolute values proves the count.
\end{proof}
Freezing $a_\tau$ to $a_0$ on $S$ performs commuting complete
source/sink moves and gives an RC MDS within two length units of
each original germ. This does not bound the difference between
distinct singular germs. Nor does \eqref{eq:mixed-wall} assert
that its partial companion exchanges are individually legal moves:
Boolean membership and acyclicity must be checked on any proposed
intermediate state.


\subsection{Ancillary atlas properties and wall examples}
\begin{proposition}[Embedding spectrum in the quadratic slice]
\label{prop:quadratic-embedding}
At $k=8$ the accidental DNA embedding ratios are
\[
 \sqrt2,\quad 2\sqrt2-1,\quad1+\sqrt2,\quad
 1+2\sqrt2,\quad3+2\sqrt2.
\]
Their total zero-word counts are respectively $288,352,648,352,576$,
including 256 structural words. At $k=2^b$, $b\ge3$, restricted
to $\rho\in\Q(\sqrt2)$, the same list is complete and
\[
 Z_{2^b}(\rho)=Z_8(\rho)^{2^{b-3}}.
\]
\end{proposition}
\begin{proof}
Antipodal pairing reduces the embedding to four entries of $D$.
The exact enumeration of $9^4$ tuples, with ordered-letter
multiplicities, checks both real and imaginary coordinates and
agrees with all $4^8$ positional words. If a zero is accidental,
at least one nonstructural coordinate is nonzero as a form, so
its ratio lies in $\Q(\sqrt2)$; no additional real ratio is omitted.
For the lift, $\{1,\zeta_{2^b},\ldots,\zeta_{2^b}^{r-1}\}$ is a
basis over $\Q(\zeta_8)$, where $r=2^{b-3}$. Separating exponents
modulo $r$ gives $r$ independent eight-position words. Scalar
and componentwise vanishing hold exactly block by block, proving
the spectrum and count. Other subfields at higher order are not
classified by this argument.
\end{proof}

\begin{proposition}[Exact family optima and a failed Galois quotient]
\label{prop:family-optima}
For the fixed selector,
\[
 \min_{\rho>1}L_6(\rho)=57,\quad
 \operatorname*{argmin}L_6=[4,\infty),
\]
\[
 \min_{\rho>1}L_8(\rho)=141,\quad
 \operatorname*{argmin}L_8=[3,1+2\sqrt2).
\]
Consequently no fixed ratio minimizes this family at every even order.
Moreover, Galois conjugation does not define an operation on terminal
equivalence classes, even when all four ratios remain admissible.
\end{proposition}
\begin{proof}
Upper ranks and attaining paths give the exact length on each
stratum. The ten optimal $k=6$ strata and sixteen optimal $k=8$
strata join to the displayed intervals; the excluded right endpoint
at $k=8$ has length 155. Since $1+2\sqrt2<4$, these loci are disjoint.

For the last assertion let
$x=5-\sqrt2$ and $y=(17+6\sqrt2)/7$.
Exact evaluation gives the same terminal at $x$ and $y$, but
their conjugates under $\sqrt2\mapsto-\sqrt2$ have residual lengths
143 and 161 and hence different terminals. All ratios exceed one.
This disproves descent to terminal classes without contradicting
the typed equality-certificate covariance.
\end{proof}

The length bound in \Cref{ext:thm:paths} is attained on all 199 strata
at $k=4,6,8$; no all-order equality is inferred. The mixed-wall
formula is checked on all 33 comparisons at eleven walls.
For example at $k=8$, $\rho=3+2\sqrt2$, the side difference is
\[
2056=8\cdot112+2\cdot540+2\cdot40,
\]
while the lengths are $(143,163,143)$. After freezing regular
midpoint changes they are $(145,163,143)$, so a width-20 excursion
persists. At $(1+2\sqrt2)/3$ the embedding-zero count remains 256
but the component changes: new zero PCRs alone cannot detect all
component changes. At $(5+4\sqrt2)/7$ the three lengths are all
149 while the components are distinct. Equal length is therefore
not a complete component invariant.


\subsection{Technical normalization details}
\label{app:component-details}

The prime-sine independence argument is proved once in \Cref{lem:prime-sines}; that proof supplies the normalization used here.

\subsubsection{Decision traces and sufficient forms}
The five-role catalogue retains a transparent branch trace. Structural forms
are registered separately from nonzero hyperplanes. The two-family theorem
(\Cref{ext:thm:two}) proves that the lambda and midpoint signs suffice;
retaining all five roles is an audit choice, not a minimality claim.
A point stratum is retained separately from its adjacent intervals.

\subsubsection{Certificate height}
Raw positional polynomials have coefficients in $\{-1,0,1\}$, with one
signed one-hot letter per position. The role-transformed Laurent
polynomials $A,B$ need not have these coefficients. For example the
midpoint polynomial of $TTT$ at $k=4$ contains the constant coefficient
$-6$, since $A(z)=\sum_{e=1}^3(z^e+z^{-e}-2)$.
For a word of length $n$, all five role templates satisfy the safe bounds
\[
 \|A\|_1,\|B\|_1\le4n,\qquad
 \|UA+dB\|_1\le4n(\|U\|_1+d)\quad(d>0).
\]
A common Laurent shift preserves this bound. The cleared degree is
$O(k+\deg U)$. These polynomial height bounds suffice for the bit
complexity statement, without asserting that every bounded polynomial
has a positional realization.

\subsection{Finite evidence and reproducibility boundary}
\label{app:finite-evidence}

The original regression proposition uses a canonical thirty-row payload.  Its frozen
identities are
\begin{align*}
  \text{producer payload prefix: }&aca5d5fc81c839c218b6eefaaf00acdc\ldots,\\
  \text{validator payload prefix: }&b79051b4190d6f13b9bb183fe1299fb6\ldots.
\end{align*}
The full SHA-256 values, schemas, exact row table, independent reconstruction
code, and replay instructions are included in the public finite-assurance
capsule distributed with this article.

\Cref{thm:twice-prime-embedding} has a separate finite arithmetic checker with payload prefix
\texttt{9ed015120913ea30}.  It verifies selected composition identities, word
counts and the $k=12$ orbit count, and uses $k=30$ only as an out-of-domain
negative control.  It is not a premise of the all-order proof.

The finite range is:
\[
  3\text{ rows at }k=4,
  \qquad 23\text{ rows at }k=6,
  \qquad 4\text{ named landmarks at }k=8.
\]
No floating-point decision enters the scientific payload.  Every row checks
the direct formula against a legal firing replay, literal RC invariance,
one-per-PCR cardinality, and acyclicity by both Kahn elimination and an
independent depth-first search.  Longest paths are counted in edges; the
corresponding vertex length is one larger for a nonempty path.

This regression capsule is exact inside its declared range. The extended
capsules supply the complete $k=8$ atlas, component classification, retiming
certificates and the 131 construction, as described in
\Cref{sec:construction-131}. Neither capsule claims a terminal atlas at
$k\ge10$, biological performance or a deduction of uniform theorems
from finite rows.
